% !TeX root = ../../Book4.tex
% 纲要标题：## 附录 D　Frobenius 范畴与稳定模
% 纲要细目（写作范围）：
% > 在正合范畴中构造稳定 Hom、合冲与余合冲自等价；完整计算双数代数的稳定模及逐项分裂复形的同伦模型，给出 Happel 三角定理的推出构造概要。

\chapter{Frobenius 范畴与稳定模}
\label{b4:app:D}
\BookChapterNumber{b4:app:D}

\begin{chapterabstract}
投射对象与内射对象相同的正合范畴，可以把经过它们的态射统一视为零。在所得稳定范畴中，合冲与余合冲互为逆操作，短正合列产生好三角。双数代数的有限维模给出一个可直接计算的例子；逐项分裂正合的复形则把这项构造具体识别为同伦范畴。
\end{chapterabstract}

\begin{learninggoals}
\begin{itemize}
\item 针对指定正合结构检验投射性、内射性及 Frobenius 条件。
\item 计算稳定 Hom，并说明合冲、余合冲的比较映射为何不依赖提升或延拓。
\item 分解双数代数的有限维模，求其稳定态射与合冲对象。
\item 把零同伦映射分解为经过可缩复形的映射，并将推出三角识别为映射锥三角。
\end{itemize}
\end{learninggoals}

令 \(A=k[\varepsilon]/(\varepsilon^2)\)。一份 \(A\)-模 \(k=A/(\varepsilon)\) 可以嵌入 \(A\)，也可以由 \(A\) 满射得到，二者放在同一列中就是
\[
0\longrightarrow k\xrightarrow{1\mapsto\varepsilon}A
\longrightarrow k\longrightarrow0.
\]
此列不分裂。如果把自由模 \(A\) 当作可以消去的部分，列的两端还保留什么联系？答案并不是把三个对象随意删去一个，而是同时改变态射的识别方式：凡经过投射模的映射都视为零。这样中项成为零对象，两个 \(k\) 却通过一个可逆的连接态射相联系。

\ProseParagraphBreak
复形中也有类似现象。映射锥 \(\operatorname{Cone}(1_X)\) 可缩，里面包含 \(X\)，余核是 \(X[1]\)。把经过可缩复形的映射视为零，恰好就是把零同伦映射视为零，所以得到同伦范畴。这两个例子共同要求一类对象既能承担投射提升，又能承担内射延拓；Frobenius 范畴把这个条件写成正合结构的性质。

\ProseParagraphBreak
本附录使用\secref{b4:sec:B01}的正合范畴、模的投射与内射性质，以及\chapref{b4:ch:06}的链同伦和映射锥。三角公理采用\cref{b4:def:triangulated-category}的约定。商范畴、移位自等价和两个模型在这里证明；一般 Frobenius 范畴的全部三角公理采用 Happel 定理，并展开其推出与八面体对象的构造，参见 Bühler\cite[Sections 12 and 13.4]{Buehler2010ExactCategories}。

% !TeX root = ../../Book4.tex
% 纲要标题：### D.1 Frobenius 范畴与稳定范畴
% 纲要细目（写作范围）：
% #### D.1.1 正合结构中的投射与内射
% #### D.1.2 把经过投射内射对象的态射视为零
% #### D.1.3 合冲与余合冲函子

\section{Frobenius 范畴与稳定范畴}
\label{b4:sec:D01}

给一个模选取投射满射后，核记录生成元之间的关系；选取内射嵌入后，余核则记录另一方向的缺失。这两个对象通常依赖选择。如果把经过投射或内射对象的映射视为零，比较映射的差别就可能消失。Frobenius 范畴使这两种消去在同一个商范畴中进行。

\subsection{正合结构中的投射与内射}
\label{b4:subsec:D01:01}

\cref{b4:def:B-exact-category}中的正合范畴只指定一部分核余核对。因此，投射性与内射性也应当针对这些可容许短正合列定义。

\begin{definition}\label{b4:def:D-relative-projective-injective}
设 \(\mathcal E\) 为局部小正合范畴。若对象 \(P\in\mathcal E\) 对每个可容许满态射 \(p:B\to C\) 都使
\[
\operatorname{Hom}_{\mathcal E}(P,B)\longrightarrow
\operatorname{Hom}_{\mathcal E}(P,C)
\]
满射，就称 \(P\) 为该正合结构中的投射对象。若对象 \(I\) 对每个可容许单态射 \(i:A\to B\) 都使
\[
\operatorname{Hom}_{\mathcal E}(B,I)\longrightarrow
\operatorname{Hom}_{\mathcal E}(A,I)
\]
满射，就称 \(I\) 为该正合结构中的内射对象。若每个对象 \(X\) 都接收一个从投射对象出发的可容许满态射，称 \(\mathcal E\) 有足够多投射对象；若每个对象都经可容许单态射嵌入一个内射对象，称它有足够多内射对象。
\end{definition}

投射性说的是所给态射可以提升，内射性说的是所给态射可以延拓。它们的有限直和以及范畴内已有的直和项仍有相同性质，因为可以逐分量提升或延拓。在交换范畴的通常正合结构中，这些定义与模论中的定义相同；在仅含分裂列的正合结构中，每个对象都投射且内射。

\begin{definition}\label{b4:def:D-frobenius-category}
设 \(\mathcal E\) 为局部小正合范畴。如果它有足够多投射对象和足够多内射对象，而且两类对象相同，就称 \(\mathcal E\) 为 \term[frobenius-fanchou]{Frobenius 范畴}[Frobenius category]。把这类共同对象组成的满子范畴记为 \(\mathcal P\)，称其中的对象为投射内射对象。
\end{definition}

同一加性范畴上的不同正合结构可能给出不同的 \(\mathcal P\)。例如，对任意域 \(k\)，在双数代数 \(A=k[\varepsilon]/(\varepsilon^2)\) 的有限维左模范畴中，\(S=A/(\varepsilon)\) 在通常正合结构下不投射，在分裂正合结构下却投射且内射，具体比较见\cref{b4:prop:D-exact-structure-stable-comparison}。下面的商构造必须连同正合结构一起理解。Frobenius 范畴的定义与基本例子见 Bühler\cite[Section 13.4]{Buehler2010ExactCategories}。

\subsection{把经过投射内射对象的态射视为零}
\label{b4:subsec:D01:02}

\begin{proposition}\label{b4:prop:D-projective-ideal}
设 \(\mathcal E\) 为局部小 Frobenius 范畴，\(\mathcal P\) 为其投射内射对象组成的满子范畴。对 \(X,Y\in\mathcal E\)，令
\[
\mathcal P(X,Y)=
\{f:X\to Y:\ f\;\text{可以分解为}\;X\to P\to Y, P\in\mathcal P\}.
\]
则 \(\mathcal P(X,Y)\) 是 \(\operatorname{Hom}_{\mathcal E}(X,Y)\) 的子群，而且与任意态射在左右两侧复合后仍属于相应子群。
\end{proposition}
\symidx{\(\mathcal P(X,Y)\)：经过投射内射对象分解的态射子群}
\begin{proof}
零映射通过零对象分解，负映射通过同一个 \(P\) 分解。若 \(f=ba\)、\(f'=b'a'\) 分别通过 \(P,P'\)，则
\[
f+f'=
\begin{pmatrix}b&b'\end{pmatrix}
\begin{pmatrix}a\\a'\end{pmatrix}
\]
通过 \(P\oplus P'\) 分解，而有限直和仍投射且内射。对 \(u:X'\to X\)、\(v:Y\to Y'\)，有 \(vfu=(vb)(au)\)，仍通过 \(P\)。
\end{proof}

\begin{definition}\label{b4:def:D-stable-category}
设 \(\mathcal E\) 为局部小 Frobenius 范畴，\(\mathcal P\) 为其投射内射对象组成的满子范畴。保留 \(\mathcal E\) 的全部对象，规定
\[
\operatorname{Hom}_{\underline{\mathcal E}}(X,Y)
=\operatorname{Hom}_{\mathcal E}(X,Y)/\mathcal P(X,Y),
\qquad [g]\circ[f]=[g\circ f],
\]
所得范畴称为 \(\mathcal E\) 的\term[wendingfanchou]{稳定范畴}[stable category]，记为 \(\underline{\mathcal E}\)。映射 \(f\) 在此商中的类称为稳定态射。
\end{definition}
\symidx{\(\underline{\mathcal E}\)：Frobenius 范畴的稳定范畴}

\cref{b4:prop:D-projective-ideal}保证复合良定义。原来的零对象和有限双积仍给出商范畴的零对象和有限双积，所以 \(\underline{\mathcal E}\) 是局部小加性范畴。这里商掉的是态射子群，既没有把对象逐个删除，也没有反演某一类态射。

\begin{proposition}\label{b4:prop:D-stable-zero}
设 \(\mathcal E\) 为局部小 Frobenius 范畴。对象 \(X\) 在 \(\underline{\mathcal E}\) 中同构于零，当且仅当 \(X\) 是投射内射对象。特别地，若 \(P\) 投射且内射，则自然包含 \(X\to X\oplus P\) 在稳定范畴中可逆。
\end{proposition}
\begin{proof}
若 \(X\in\mathcal P\)，恒等映射通过 \(X\) 分解，故 \([1_X]=0\)。反过来，\([1_X]=0\) 给出 \(1_X=ba\)，其中 \(a:X\to P\)、\(b:P\to X\)、\(P\in\mathcal P\)。给定可容许满态射 \(q:E\to Y\) 及 \(f:X\to Y\)，先由 \(P\) 的投射性把 \(fb\) 提升为 \(\widetilde f:P\to E\)，满足 \(q\widetilde f=fb\)。于是
\[
 q(\widetilde f a)=fba=f,
\]
所以 \(\widetilde f a\) 提升 \(f\)，证明 \(X\) 投射。对可容许单态射 \(j:Y\to Z\) 及 \(g:Y\to X\)，由 \(P\) 的内射性把 \(ag\) 延拓为 \(\widetilde g:Z\to P\)，满足 \(\widetilde g j=ag\)。此时
\[
 (b\widetilde g)j=bag=g,
\]
证明 \(X\) 内射。这里已经有实际的收缩数据 \(a,b\)，不要求 \(P\) 的互补对象存在，也不要求范畴幂等完备。最后，包含与投影的一个复合是恒等，另一个与恒等之差只作用于 \(P\) 分量，故通过 \(P\) 分解。
\end{proof}

\subsection{合冲与余合冲函子}
\label{b4:subsec:D01:03}

在 Frobenius 范畴中，为每个对象固定可容许短正合列
\begin{equation}\label{b4:eq:D-syzygy-sequences}
0\longrightarrow\Omega X\xrightarrow{j_X}P_X\xrightarrow{p_X}X\longrightarrow0,
\qquad
0\longrightarrow X\xrightarrow{i_X}I_X\xrightarrow{q_X}\Sigma X\longrightarrow0,
\end{equation}
其中 \(P_X,I_X\in\mathcal P\)。称 \(\Omega X\) 为这次选择的合冲对象，\(\Sigma X\) 为余合冲对象。对大范畴，这些逐对象指定作为已选数据，沿用\subsecref{b4:subsec:0101:03}的大小与选择约定。
\symidx{\(\Omega X,\ \Sigma X\)：稳定范畴中的合冲与余合冲对象}

\begin{theorem}\label{b4:thm:D-stable-syzygy-equivalence}
设 \(\mathcal E\) 为局部小 Frobenius 范畴，并固定式\eqref{b4:eq:D-syzygy-sequences}中的序列。核与余核的比较映射分别定义加性函子
\[
\Omega,\Sigma:\underline{\mathcal E}\longrightarrow\underline{\mathcal E}.
\]
更换已选序列得到的函子与原函子自然同构，而且
\[
\Omega\Sigma\cong\operatorname{Id}_{\underline{\mathcal E}},
\qquad
\Sigma\Omega\cong\operatorname{Id}_{\underline{\mathcal E}}.
\]
因此 \(\Sigma\) 是自等价，\(\Omega\) 是其拟逆。
\end{theorem}
\begin{proof}
给定 \(f:X\to Y\)，投射性提供 \(\widetilde f:P_X\to P_Y\)，满足 \(p_Y\widetilde f=fp_X\)。固定这张提升后，核的泛性质才给出唯一 \(a:\Omega X\to\Omega Y\)，使 \(j_Ya=\widetilde f j_X\)；提升本身通常不唯一。若换成另一提升 \(\widetilde f'\)，两提升之差唯一写成 \(j_Yc\)，其中 \(c:P_X\to\Omega Y\)。由 \(j_Y(a-a')=j_Ycj_X\) 及 \(j_Y\) 为单态射，得到 \(a-a'=cj_X\)。这个差通过 \(P_X\) 分解，故 \([a]=[a']\)。构造依赖提升的选择，但商去这项误差后，稳定类不再依赖它。

若 \(f\) 通过投射对象 \(Q\) 分解为 \(X\xrightarrow{u}Q\xrightarrow{v}Y\)，先把 \(v\) 提升为 \(\widetilde v:Q\to P_Y\)。取 \(\widetilde f=\widetilde vup_X\)，则 \(\widetilde f j_X=0\)，故所诱导的核映射为零。这说明构造只依赖 \([f]\)。两提升的复合是复合态射的一种提升，提升之和是态射之和的一种提升，恒等可以用恒等提升；由刚证的独立性，得到加性函子 \(\Omega\)。

内射性同样提供 \(\widehat f:I_X\to I_Y\)，满足 \(\widehat f i_X=i_Yf\)。余核的泛性质给出 \(b:\Sigma X\to\Sigma Y\)。两个延拓之差写成 \(c q_X\)，其中 \(c:\Sigma X\to I_Y\)，所以两个余核映射之差为 \(q_Yc\)，通过 \(I_Y\) 分解。若 \(f\) 通过内射对象 \(Q\) 分解，先把 \(X\to Q\) 沿 \(i_X\) 延拓，再接 \(Q\to Y\to I_Y\)，所得余核映射为零。复合、加法与恒等的相容性由同样的比较得到，因此 \(\Sigma\) 也是加性函子。

现在比较同一对象 \(X\) 的两套投射序列，记其满态射为 \(p_X:P_X\to X\) 与 \(p'_X:P'_X\to X\)。投射性给出 \(u:P_X\to P'_X\)、\(v:P'_X\to P_X\)，满足 \(p'_Xu=p_X\)、\(p_Xv=p'_X\)，从而诱导两套核之间的比较映射。其复合由 \(vu\) 诱导，而 \(p_Xvu=p_X\)，所以 \(vu\) 与 \(1_{P_X}\) 都是 \(1_X\) 的提升。前面已证明两个提升所诱导的核映射之差通过 \(P_X\) 分解；恒等提升诱导恒等核映射，因此该复合在稳定范畴中为恒等，反向同理。对 \(f:X\to Y\)，两条比较路径都来自 \(f\) 的提升；其差通过相应投射对象分解，故这些比较同构自然。内射序列的比较由上述延拓论证得到相同结论。

第二条序列原来是 \(X\) 的内射序列；因 \(I_X\) 同时投射，它又可作 \(\Sigma X\) 的投射序列，其核就是 \(X\)。构造 \(\Sigma[f]\) 时的延拓 \(\widehat f:I_X\to I_Y\) 满足 \(q_Y\widehat f=(\Sigma f)q_X\)，所以也可作 \(\Sigma f\) 的投射提升。它又满足 \(\widehat f i_X=i_Yf\)，而 \(i_Y\) 为单态射，故限制在核 \(X\) 上所诱导的映射恰为 \(f\)。因此用这组序列计算 \(\Omega\Sigma\) 得到恒等函子，再与固定的投射序列比较，就得到自然同构 \(\Omega\Sigma\cong\operatorname{Id}\)。

第一条序列原来是 \(X\) 的投射序列；因 \(P_X\) 同时内射，它又可作 \(\Omega X\) 的内射序列，其余核就是 \(X\)。对 \(f\) 使用构造 \(\Omega[f]\) 时的提升图，余核上的映射也恰为 \(f\)。用这组序列计算 \(\Sigma\Omega\) 因而得到恒等函子，再与固定的内射序列比较，得到自然同构 \(\Sigma\Omega\cong\operatorname{Id}\)。两种同构的自然性都来自刚才的序列比较，而不只是在每个对象上单独得到同构。
\end{proof}

核和余核在原范畴中可以不同；自然同构发生在稳定范畴中。投射内射项同时能消去两种比较误差，正是合冲与余合冲成为互逆操作的原因。

\begin{proposition}\label{b4:prop:D-exact-structure-stable-comparison}
设 \(k\) 为域，\(A=k[\varepsilon]/(\varepsilon^2)\)，\(S=A/(\varepsilon)\)，并在有限维幺性左 \(A\)-模的同一个加性范畴上分别指定通常短正合结构和分裂短正合结构。两者都是 Frobenius 正合结构，但所得稳定范畴不同。

在通常结构中，投射内射对象恰为有限自由模 \(A^r\)，且
\[
 0\longrightarrow S\xrightarrow{\iota}A\xrightarrow{\pi}S
 \longrightarrow0,\qquad\iota(\overline1)=\varepsilon,
\]
是不分裂的短正合列，其中 \(\pi\) 为商映射。此时
\[
 \operatorname{End}_{\underline{A\text{-}\mathrm{mod}}}(S)\cong k,
\]
所以 \(S\) 在稳定范畴中不为零。在分裂结构中，每个对象都投射且内射，稳定范畴等价于零范畴。
\end{proposition}
\begin{proof}
\subsecref{b4:subsec:D02:02}用取 \(\varepsilon\) 系数的泛函证明左模 \(A\cong D(A)\)，由\cref{b4:prop:D-frobenius-module-category}，通常结构为 Frobenius 结构。\cref{b4:prop:D-dual-number-stable}给出全部有限维模的分解 \(A^r\oplus S^s\)，并证明投射内射对象恰为 \(A^r\)。该证明还说明 \(\pi:A\to S\) 没有 \(A\)-线性截面，所以所列短正合列不分裂。任何 \(S\to A^r\) 的像落在 \(\varepsilon A^r\)，任何 \(A^r\to S\) 又杀掉它，因而经过投射内射对象的 \(S\to S\) 映射全为零。因此稳定自同态环仍是 \(\operatorname{End}_A(S)\cong k\)，其恒等类非零。

只含分裂列的结构由\cref{b4:prop:B-split-exact}给出。对分裂满态射 \(q:B\to C\)，取截面 \(s:C\to B\)，则任意 \(f:X\to C\) 都由 \(sf\) 提升。对分裂单态射 \(j:B\to C\)，取收缩 \(r:C\to B\)，则任意 \(g:B\to X\) 都由 \(gr\) 延拓。所以每个对象都投射且内射，恒等态射本身提供足够多投射与内射对象。每个态射 \(f:X\to Y\) 都通过投射内射对象 \(X\) 分解，稳定 Hom 群全部为零，所得范畴遂等价于零范畴。
\end{proof}

% !TeX root = ../../Book4.tex
% 纲要标题：### D.2 稳定模与同伦范畴
% 纲要细目（写作范围）：
% #### D.2.1 Frobenius 代数的有限维模
% #### D.2.2 双数代数的稳定模
% #### D.2.3 逐项分裂正合的复形

\section{稳定模与同伦范畴}
\label{b4:sec:D02}

稳定范畴的定义能用于相当不同的对象。一个模型从有限维代数出发：某些代数的投射模本身也内射。另一个模型从复形出发：若只指定逐项分裂的短正合列，可缩复形恰好承担投射与内射的共同角色。前者让我们计算一个具体模的稳定态射，后者则重新得到本卷已经熟悉的同伦范畴。

\subsection{Frobenius 代数的有限维模}
\label{b4:subsec:D02:01}

设 \(k\) 为域，\(A\) 为有限维含幺结合 \(k\)-代数，未必交换。记 \(A\text{-}\mathrm{mod}\) 为有限维幺性左 \(A\)-模及其 \(A\)-线性映射组成的范畴，配备通常的短正合结构。对偶空间
\[
D(A)=\operatorname{Hom}_k(A,k)
\]
有左 \(A\)-作用 \((a\varphi)(b)=\varphi(ba)\)。这里出现 \(ba\) 而非 \(ab\)，是为使所定义的作用成为左作用。
\symidx{\(A\text{-}\mathrm{mod}\)：有限维代数上的有限维左模范畴}

\begin{definition}\label{b4:def:D-frobenius-algebra}
设 \(k\) 为域，\(A\) 为有限维含幺结合 \(k\)-代数。在 \(D(A)=\operatorname{Hom}_k(A,k)\) 上指定左作用 \((a\varphi)(b)=\varphi(ba)\)。如果左正则模 \({}_AA\) 与 \({}_AD(A)\) 同构，就称 \(A\) 为 \term[frobenius-daishu]{Frobenius 代数}[Frobenius algebra]。
\end{definition}
\symidx{\(D(A)\)：有限维代数的线性对偶及其左模结构}

\begin{proposition}\label{b4:prop:D-frobenius-module-category}
设 \(k\) 为域，\(A\) 为 \(k\) 上的有限维 Frobenius 代数。则通常正合结构下的 \(A\text{-}\mathrm{mod}\) 是 Frobenius 范畴，投射内射对象恰为有限生成投射左 \(A\)-模。
\end{proposition}
\begin{proof}
对每个有限维左 \(A\)-模 \(M\)，有自然同构
\[
\operatorname{Hom}_A(M,D(A))\longrightarrow D(M),
\qquad f\longmapsto\bigl(m\mapsto f(m)(1)\bigr).
\]
其逆把 \(\lambda\in D(M)\) 送到
\[
f_\lambda(m)(b)=\lambda(bm).
\]
确有 \(f_\lambda(am)(b)=\lambda(bam)=(a f_\lambda(m))(b)\)，故逆映射左 \(A\)-线性。若 \(f\) 左线性，则
\(f(m)(b)=(b f(m))(1)=f(bm)(1)\)，所以两映射互逆。有限维向量空间的对偶函子正合，因此 \(D(A)\) 内射。假设 \(A\cong D(A)\) 便使左正则模 \(A\) 内射，有限直和及其直和项也内射；每个有限生成投射模是某个 \(A^r\) 的直和项，因而内射。

每个 \(M\) 接收有限自由模的满射，所以有足够多投射对象。另取 \(D(M)\) 的一组基 \(\lambda_1,\ldots,\lambda_d\)，以上公式给出单射
\[
M\longrightarrow D(A)^d,
\qquad m\longmapsto(f_{\lambda_1}(m),\ldots,f_{\lambda_d}(m)).
\]
若像为零，在 \(b=1\) 处求值便得所有 \(\lambda_j(m)=0\)，故 \(m=0\)。这证明有足够多内射对象。若 \(M\) 本身内射，上述单射分裂，于是 \(M\) 是 \(D(A)^d\cong A^d\) 的直和项，因而投射。两类对象相同，结论成立。
\end{proof}

这个证明只使用左模同构 \(A\cong D(A)\)，没有把 \(A\) 当作交换代数。稳定范畴 \(\underline{A\text{-}\mathrm{mod}}\) 因此把经过有限生成投射左模的态射视为零。

\subsection{双数代数的稳定模}
\label{b4:subsec:D02:02}

取任意域 \(k\) 上的双数代数
\[
A=k[\varepsilon]/(\varepsilon^2).
\]
令 \(\lambda(a+b\varepsilon)=b\)。映射
\[
\phi:A\longrightarrow D(A),\qquad
\phi(x)(y)=\lambda(yx)
\]
左 \(A\)-线性，因为 \(\phi(ax)(y)=\lambda(yax)=(a\phi(x))(y)\)。若 \(x=a+b\varepsilon\ne0\)，则 \(b\ne0\) 时 \(\lambda(x)=b\ne0\)，\(a\ne0\) 时 \(\lambda(\varepsilon x)=a\ne0\)。故 \(\phi\) 单射；两边维数均为二，所以同构。由\cref{b4:prop:D-frobenius-module-category}，\(A\text{-}\mathrm{mod}\) 是 Frobenius 范畴。

\begin{proposition}\label{b4:prop:D-dual-number-stable}
设 \(k\) 为域，\(A=k[\varepsilon]/(\varepsilon^2)\)，并把 \(k=A/(\varepsilon)\) 视为左 \(A\)-模。每个有限维左 \(A\)-模都同构于
\[
A^r\oplus k^s\qquad(r,s\ge0).
\]
其中投射模恰为 \(A^r\)，并且
\[
\operatorname{Hom}_{\underline{A\text{-}\mathrm{mod}}}(k^s,k^t)
\cong M_{t\times s}(k),
\qquad
\Omega k\cong k\cong\Sigma k.
\]
把有限维 \(k\)-向量空间赋予零 \(\varepsilon\) 作用，给出到 \(\underline{A\text{-}\mathrm{mod}}\) 的加性等价。
\end{proposition}
\begin{proof}
一个 \(A\)-模就是带有平方为零算子 \(T(m)=\varepsilon m\) 的向量空间。取 \(\operatorname{im}T\) 的基 \(u_1,\ldots,u_r\)，选 \(v_i\) 使 \(Tv_i=u_i\)，再把 \(u_i\) 扩充为 \(\ker T\) 的基 \(u_i,w_1,\ldots,w_s\)。这些 \(u_i,v_i,w_j\) 构成全空间的基：应用 \(T\) 可先消去任何线性关系中的 \(v_i\) 系数；每个向量减去适当的 \(v_i\) 组合后则落入 \(\ker T\)。每对 \(v_i,u_i\) 给出一份 \(A\)，每个 \(w_j\) 给出一份 \(k\)，得到分解。

模 \(k\) 不投射。若商映射 \(\pi:A\to k\) 有 \(A\)-线性截面 \(s\)，则 \(\varepsilon s(1)=s(\varepsilon\cdot1)=0\)，故 \(s(1)\in\varepsilon A\)，于是 \(\pi s(1)=0\)，与截面条件矛盾。因此含有 \(k\) 直和项的模不投射，而 \(A^r\) 自由，投射模恰为这些对象。

任意 \(A\)-线性映射 \(k^s\to A^r\) 的像落入 \(\varepsilon A^r\)，而任何 \(A^r\to k^t\) 都把该子模送到零。因此两者的复合为零。所有经过投射模的 \(k^s\to k^t\) 映射均为零，稳定 Hom 就等于普通 Hom，即所给矩阵空间。每个 \(A^r\oplus k^s\) 在稳定范畴中同构于 \(k^s\)，故向量空间的这个函子全忠实且本质满。具体的拟逆可取
\[
M\longmapsto\ker(\varepsilon:M\to M)/\varepsilon M;
\]
它把 \(A\) 送到零，把 \(k\) 送到 \(k\)，并由上面的分解给出同一稳定 Hom 识别。

最后，短正合列
\begin{equation}\label{b4:eq:D-dual-number-extension}
0\longrightarrow k\xrightarrow{\iota}A\xrightarrow{\pi}k\longrightarrow0,
\qquad \iota(1)=\varepsilon,
\end{equation}
既可作右端 \(k\) 的投射序列，也可作左端 \(k\) 的内射序列。于是 \(\Omega k=\varepsilon A\cong k\)，\(\Sigma k=A/\varepsilon A\cong k\)。标量映射可提升、延拓为 \(A\) 上的同一标量映射，所以这些识别也保持态射。在向量空间模型中，\(\Omega\) 与 \(\Sigma\) 因而都自然同构于恒等函子。
\end{proof}

上述等价是加性范畴的等价。原来的模范畴有非分裂扩张\eqref{b4:eq:D-dual-number-extension}，稳定商却把中项 \(A\) 变成零。在 \(\Sigma k\cong k\) 的选择下，这条扩张给出的好三角为
\[
k\longrightarrow0\longrightarrow k\xrightarrow{-1_k}k,
\]
连接态射的计算见\cref{b4:prop:D-conflation-triangle}后的双数代数例子。这并不是有限维向量空间范畴中的短正合列；加性等价本身并未将原来的模短正合结构搬到向量空间上。

\ProseParagraphBreak
双数代数中 \(\Omega\cong\operatorname{Id}\) 也不是 Frobenius 条件的一般后果。一般结论是 \(\Omega\) 与 \(\Sigma\) 互为拟逆；对三次截断多项式环的简单模，二次合冲才回到原模，一次合冲并不与原模稳定同构，具体计算见\cref{b4:prop:D-truncated-cubic-period}。

\subsection{逐项分裂正合的复形}
\label{b4:subsec:D02:03}

设 \(\mathcal A\) 为局部小加性范畴。上链复形与复形映射的定义仍可使用，因为只需态射加法、复合及零态射。给 \(\operatorname{Ch}(\mathcal A)\) 指定在每个次数都分裂的短正合列。分裂只针对各项，不要求所选截面组成复形映射。

\ProseParagraphBreak
例如，设 \(k\) 为域，\(D=(k\xrightarrow{1_k}k)\) 位于次数 \(0,1\)，\(S^i\) 是只在次数 \(i\) 取值 \(k\) 的复形。短正合列
\[
0\longrightarrow S^1\longrightarrow D\longrightarrow S^0\longrightarrow0
\]
逐项分裂，却不在复形范畴中分裂。若商映射有复形截面 \(s\)，零次截面条件要求 \(s^0=1_k\)，复形映射条件却要求
\[
\mathrm{d}_D^0s^0=s^1\mathrm{d}_{S^0}^0=0,
\]
而左边为 \(1_k\)。每度存在截面，只解决分次对象的分裂，尚未解决与微分的相容性。

\begin{proposition}\label{b4:prop:D-complex-frobenius}
设 \(\mathcal A\) 为局部小加性范畴。逐项分裂的短正合列给出 \(\operatorname{Ch}(\mathcal A)\) 的正合结构；在此结构下它是 Frobenius 范畴，投射内射对象恰为可缩复形。
\end{proposition}
\begin{proof}
对任意复形 \(X,Y\)，复形映射 \(X\to Y\) 由满足 \(\mathrm{d}_Y^nf^n=f^{n+1}\mathrm{d}_X^n\) 的态射族 \((f^n)_{n\in\mathbb Z}\) 给出。这些态射族构成逐项 Hom 集的可数积
\[
\prod_{n\in\mathbb Z}\operatorname{Hom}_{\mathcal A}(X^n,Y^n)
\]
的子集。因为 \(\mathcal A\) 局部小，这个可数积是集合，故 \(\operatorname{Ch}(\mathcal A)\) 也局部小。

逐项分裂列的核与余核在每个次数都有泛性质，所得态射族与微分相容，所以也在复形范畴中互为核与余核。选逐项分裂后，中项的微分有形式
\[
Y^n=X^n\oplus Z^n,
\qquad
\mathrm{d}_Y^n=
\begin{pmatrix}\mathrm{d}_X^n&t^n\\0&\mathrm{d}_Z^n\end{pmatrix},
\quad
\mathrm{d}_X^{n+1}t^n+t^{n+1}\mathrm{d}_Z^n=0.
\]
右上角 \(t^n\) 记录的正是逐项分裂后微分仍可能把商的坐标送入核的坐标；上面的两项复形例子中，它就是 \(1_k\)。
沿复形映射 \(f:X\to X'\) 推出，取 \(Y'^n=X'^n\oplus Z^n\)，把右上角改为 \(f^{n+1}t^n\)；这满足复形条件，逐项推出的泛性质也保证所得复形有推出泛性质。沿 \(g:Z'\to Z\) 拉回，则取 \(X^n\oplus Z'^n\)，右上角改为 \(t^ng^n\)。两个逐项分裂包含的复合逐项仍是分裂包含，商由两个互补项的直和逐项给出；两个逐项分裂投影的复合及其核有对偶构造。因此全部正合结构公理成立。

若 \(K\) 可缩，固定收缩 \(h_K\)，满足 \(\mathrm{d}_Kh_K+h_K\mathrm{d}_K=1_K\)。对任意复形 \(X\)，Hom 复形 \(\operatorname{Hom}^\bullet(K,X)\) 与 \(\operatorname{Hom}^\bullet(X,K)\) 都可缩。具体地，对次数 \(n\) 的齐次映射 \(a\)，两者的收缩分别为
\[
a\longmapsto(-1)^n a h_K,
\qquad
a\longmapsto h_Ka.
\]
代入 \(\partial a=\mathrm{d}a-(-1)^na\mathrm{d}\)，两次交叉项相消，得到 \(\partial H+H\partial=1\)。

给定逐项分裂满射 \(p:E\to Y\) 及复形映射 \(f:K\to Y\)，先逐项提升为次数零映射 \(a_0:K\to E\)，满足 \(pa_0=f\)。逐项提升未必与微分相容，其误差为 \(\partial a_0\)。因为 \(p\partial a_0=\partial f=0\)，误差经核 \(i:X\to E\) 唯一写成
\[
\partial a_0=i r,\qquad r\in\operatorname{Hom}^1(K,X).
\]
又有 \(i\partial r=\partial^2a_0=0\)，而 \(i\) 逐项单射，所以 \(\partial r=0\)。核复形中的误差因此是一个一次余圈。Hom 复形的收缩把它变成余边：令 \(b=H(r)\)，收缩等式给出 \(\partial b=r\)。修正逐项提升为
\[
a=a_0-i b,\qquad
\partial a=ir-i\partial b=0,\qquad pa=f.
\]
于是 \(a\) 是所需的复形映射，证明 \(K\) 投射。这里不需要 \(K\) 的各项在 \(\mathcal A\) 中投射；逐项分裂提供初始提升，整个可缩复形的同伦则消去其误差。

对偶地，给定逐项分裂单射 \(i:X\to E\) 及复形映射 \(f:X\to K\)，先逐项延拓为 \(a_0:E\to K\)。其误差满足 \((\partial a_0)i=0\)，所以经余核 \(q:E\to Y\) 写成 \(rq\)，其中 \(r\in\operatorname{Hom}^1(Y,K)\)。由 \(\partial^2a_0=0\) 及 \(q\) 逐项满射，仍有 \(\partial r=0\)。用 \(\operatorname{Hom}^\bullet(Y,K)\) 的收缩取 \(b=H(r)\)，则 \(a_0-bq\) 的微分为零，且在 \(X\) 上仍等于 \(f\)。这给出复形映射延拓，证明 \(K\) 内射。

对任意 \(X\)，令
\[
I_X=\operatorname{Cone}(1_X),\qquad
I_X^n=X^n\oplus X^{n+1},\qquad
\mathrm{d}_{I_X}^n=
\begin{pmatrix}\mathrm{d}_X^n&1\\0&-\mathrm{d}_X^{n+1}\end{pmatrix}.
\]
映射 \(h^n(y,x)=(0,y)\) 是收缩。有逐项分裂短正合列
\begin{equation}\label{b4:eq:D-contractible-embedding}
0\longrightarrow X\xrightarrow{x\mapsto(x,0)}I_X
\xrightarrow{(y,x)\mapsto x}X[1]\longrightarrow0.
\end{equation}
这给出足够多内射对象。将序列移位 \(-1\)，得到逐项分裂短正合列
\[
0\longrightarrow X[-1]\longrightarrow I_X[-1]\longrightarrow X\longrightarrow0.
\]
\(I_X[-1]\) 仍可缩，因而由已证部分投射；到 \(X\) 的可容许满射给出足够多投射对象。任何内射对象都是相应 \(I_X\) 的直和项，任何投射对象都是相应 \(I_X[-1]\) 的直和项；若 \(ba=1_X\) 且 \(K\) 的收缩为 \(h\)，则 \(bha\) 收缩 \(X\)。因此这些直和项都可缩，两类对象恰为可缩复形。
\end{proof}

\begin{theorem}\label{b4:thm:D-stable-complex-homotopy}
设 \(\mathcal A\) 为局部小加性范畴，\(\operatorname{Ch}(\mathcal A)\) 配备逐项分裂正合结构。复形映射 \(f:X\to Y\) 通过可缩复形分解，当且仅当 \(f\) 零同伦。因此恒等对象对应诱导加性同构
\[
\underline{\operatorname{Ch}(\mathcal A)}\cong K(\mathcal A).
\]
在式\eqref{b4:eq:D-contractible-embedding}的选择下，稳定范畴的 \(\Sigma\) 就是复形移位 \([1]\)。
\end{theorem}
\begin{proof}
若 \(f=ba\) 通过可缩复形 \(K\)，取其收缩 \(h_K\)，则
\[
f=b(\mathrm{d}_Kh_K+h_K\mathrm{d}_K)a
=\mathrm{d}_Y(bh_Ka)+(bh_Ka)\mathrm{d}_X,
\]
所以零同伦。反过来，若 \(f=\mathrm{d}_Yh+h\mathrm{d}_X\)，定义
\[
b:I_X\longrightarrow Y,
\qquad b^n(y,x)=f^n(y)+h^{n+1}(x).
\]
复形条件在第二个分量上恰好为
\(\mathrm{d}_Y^nh^{n+1}=f^{n+1}-h^{n+2}\mathrm{d}_X^{n+1}\)，由零同伦等式成立；第一个分量则由 \(f\) 为复形映射保证。于是 \(b\) 是复形映射，并且 \(f=b\circ i_X\)。这给出经过可缩 \(I_X\) 的分解。

两个复形映射的稳定类相同，等价于它们之差零同伦，故商 Hom 群与 \(K(\mathcal A)\) 的 Hom 群完全相同，且保持复合。序列\eqref{b4:eq:D-contractible-embedding}的余核是 \(X[1]\)；对 \(f\) 的延拓可取 \((y,x)\mapsto(fy,fx)\)，诱导的余核映射正是 \(f[1]\)。因此移位也相同。
\end{proof}

这个识别使用逐项分裂的正合结构，对任意局部小加性范畴 \(\mathcal A\) 都成立，也不要求它幂等完备。若 \(\mathcal A\) 进一步是交换范畴，则可讨论零调复形与拟同构。在 \(K(\mathcal A)\) 中将拟同构变成同构，等价于把零调复形变成零，所得导出范畴为
\[
D(\mathcal A)=K(\mathcal A)[\mathrm{qis}^{-1}],
\]
其中 \(\mathrm{qis}\) 表示拟同构组成的态射类。稳定商消去可缩复形，导出局部化还会消去不可缩的零调复形。例如，对域 \(k\) 上的 \(A=k[\varepsilon]/(\varepsilon^2)\)，令 \(E\) 每个整数次数都为 \(A\)，每个微分均为乘 \(\varepsilon\)。\cref{b4:exm:unbounded-free-not-k-projective}已证明它零调但不可缩。因此 \(E\) 在同伦范畴及相应的逐项分裂稳定范畴中不为零，而 \(E\to0\) 是拟同构，故它在导出范畴中为零。复形的这一 Frobenius 模型也见 Bühler\cite[Section 13.4, Example 13.8]{Buehler2010ExactCategories}。

\begin{proposition}\label{b4:prop:D-truncated-cubic-period}
设 \(k\) 为域，\(A=k[t]/(t^3)\)，\(S=A/(t)\)，并在有限维幺性左 \(A\)-模范畴中取通常短正合结构。则 \(A\) 为 Frobenius 代数。在由商映射 \(A\to S\) 及乘 \(t\) 的满射 \(A\to tA\) 所作的合冲选择下，
\[
\Omega S=tA,\qquad \Omega^2S=t^2A\cong S,
\qquad
\operatorname{End}_{\underline{A\text{-}\mathrm{mod}}}(S)\cong k.
\]
但 \(\Omega S\) 与 \(S\) 在稳定范畴中不同构，所以 \(S\) 的合冲周期恰为二。
\end{proposition}
\begin{proof}
令 \(\lambda:A\to k\) 取 \(t^2\) 的系数，定义
\[
\phi:A\longrightarrow D(A),\qquad \phi(a)(b)=\lambda(ba).
\]
按 \(D(A)\) 的左作用，有 \(\phi(ca)(b)=\lambda(bca)=(c\phi(a))(b)\)，所以 \(\phi\) 左 \(A\)-线性。若非零 \(a\) 的最低非零次数为 \(j\in\{0,1,2\}\)，取 \(b=t^{2-j}\)，则 \(\lambda(ba)\ne0\)，故 \(\phi\) 单射。两边的 \(k\) 维数都是三，所以它同构，证明 Frobenius 性。由\cref{b4:prop:D-frobenius-module-category}，模范畴因此是 Frobenius 范畴，有限生成投射模正是所消去的投射内射对象。

商映射 \(A\to S\) 的核为 \(tA\)，而乘 \(t\) 的满射 \(A\to tA\) 的核为 \(t^2A\)。两张满射的来源都是投射模，故给出所列合冲。映射 \(S\to t^2A\)、\(\overline1\mapsto t^2\) 是左模同构，所以二次合冲回到 \(S\)。

普通自同态环 \(\operatorname{End}_A(S)\) 为 \(k\)。任意 \(S\to A^r\) 的像被 \(t\) 零化，因而落入 \(t^2A^r\)，而任何 \(A^r\to S\) 都杀掉该子模。所以经过有限自由模的 \(S\to S\) 映射为零。有限生成投射模是某个 \(A^r\) 的直和项，经过它的分解也可写成经过 \(A^r\) 的分解。因此商去投射分解后，自同态环仍为 \(k\)，特别有 \([1_S]\ne0\)。

任意 \(u:S\to tA\) 的像都被 \(t\) 零化，所以包含于 \(t^2A\)。任意 \(v:tA\to S\) 又满足 \(v(t^2)=t\cdot v(t)=0\)，故 \(vu=0\)。若 \(tA\) 与 \(S\) 稳定同构，就能选取这样两张映射，使其复合的稳定类为 \([1_S]\)，与刚证的 \([1_S]\ne0\) 矛盾。因此一次合冲未回到原对象，而二次合冲已回到它。
\end{proof}

% !TeX root = ../../Book4.tex
% 纲要标题：### D.3 短正合列产生的三角
% 纲要细目（写作范围）：
% ---
% #### D.3.1 推出构造与连接态射
% #### D.3.2 三角结构及映射锥

\section{短正合列产生的三角}
\label{b4:sec:D03}

把投射内射对象变成零以后，短正合列的中项和两端之间仍有联系。例如双数代数的列 \(0\to k\to A\to k\to0\) 不能只剩下两个零箭头：它的两个端项由余合冲操作联系起来。下面用推出构造保留这项连接资料，并与复形的映射锥比较。

\subsection{推出构造与连接态射}
\label{b4:subsec:D03:01}

固定 Frobenius 范畴中的序列 \(0\to X\xrightarrow{i_X}I_X\xrightarrow{q_X}\Sigma X\to0\)。对任意 \(f:X\to Y\)，沿 \(f\) 推出可容许单态射 \(i_X\)，得到
\begin{equation}\label{b4:eq:D-standard-pushout}
\begin{tikzcd}[column sep=large,row sep=large]
X\arrow[r,"i_X"]\arrow[d,"f"']&I_X\arrow[r,"q_X"]\arrow[d,"a_f"]&\Sigma X\arrow[d,equal]\\
Y\arrow[r,"b_f"']&C_f\arrow[r,"c_f"']&\Sigma X.
\end{tikzcd}
\end{equation}
下行是可容许短正合列。推出也给出可容许短正合列
\begin{equation}\label{b4:eq:D-graph-inflation}
0\longrightarrow X\xrightarrow{(f,-i_X)}Y\oplus I_X
\xrightarrow{(b_f,a_f)}C_f\longrightarrow0.
\end{equation}
这里第一个映射是可容许单态射：它由分裂包含 \(X\to X\oplus Y\)、变换 \((x,y)\mapsto(x,y+fx)\)、可容许单态射 \((-i_X)\oplus1_Y\) 及交换两项的同构复合而成。推出的泛性质给出第二个映射的余核性质，核与可容许性由正合结构的推出性质保证；其一般表述见 Bühler\cite[Proposition 2.12]{Buehler2010ExactCategories}。

\begin{definition}\label{b4:def:D-stable-standard-triangle}
设 \(\mathcal E\) 为局部小 Frobenius 范畴，固定余合冲序列并由式\eqref{b4:eq:D-standard-pushout}构造 \(C_f\)。在 \(\underline{\mathcal E}\) 中，把
\[
X\xrightarrow{[f]}Y\xrightarrow{[b_f]}C_f\xrightarrow{[c_f]}\Sigma X
\]
称为标准三角。与某个标准三角同构的三角称为好三角。
\end{definition}

若 \(f'\) 与 \(f\) 之差通过 \(P\in\mathcal P\) 分解为 \(vu\)，可把 \(u:X\to P\) 延拓为 \(t:I_X\to P\)。推出的泛性质给出 \(C_f\to C_{f'}\)：在 \(Y\) 上用恒等，在 \(I_X\) 上用 \(a_{f'}-b_{f'}vt\)。反向把负号改为正号，两复合在 \(Y,I_X\) 上都是恒等，因而互逆；它们也保持到 \(\Sigma X\) 的映射。因此标准三角的同构类型只依赖稳定态射。

\begin{proposition}\label{b4:prop:D-conflation-triangle}
设 \(\mathcal E\) 为局部小 Frobenius 范畴，固定余合冲序列。给定可容许短正合列
\[
0\longrightarrow X\xrightarrow{u}Y\xrightarrow{v}Z\longrightarrow0,
\]
把 \(i_X:X\to I_X\) 沿 \(u\) 延拓为 \(a:Y\to I_X\)，并由余核定义 \(\delta:Z\to\Sigma X\)，满足
\begin{equation}\label{b4:eq:D-connecting-sign}
\delta v=-q_Xa.
\end{equation}
则 \([\delta]\) 不依赖延拓的选择，而且
\[
X\xrightarrow{[u]}Y\xrightarrow{[v]}Z\xrightarrow{[\delta]}\Sigma X
\]
是好三角。反过来，每个好三角都同构于某个可容许短正合列按此方式给出的三角。
\end{proposition}
\begin{proof}
由于 \(q_Xau=q_Xi_X=0\)，\(\delta\) 存在且唯一。两个延拓之差经 \(v\) 写成 \(dv\)，其中 \(d:Z\to I_X\)，所以两个连接态射之差为 \(-q_Xd\)，通过 \(I_X\) 分解。

考虑 \(u\) 的推出 \(C_u\)。把原短正合列沿 \(i_X\) 推出，还得到可容许短正合列
\[
0\longrightarrow I_X\xrightarrow{a_u}C_u\xrightarrow{r}Z\longrightarrow0,
\qquad rb_u=v.
\]
推出泛性质由 \(1_{I_X}\) 和 \(a:Y\to I_X\) 给出 \(s:C_u\to I_X\)，满足 \(sa_u=1\)、\(sb_u=a\)。所以此列分裂，\(C_u\cong I_X\oplus Z\)，\([r]\) 在稳定范畴中可逆。在推出的两个来源上分别检查，有
\[
c_u=q_Xs+\delta r:
\quad q_Xsa_u+\delta ra_u=q_X,
\quad q_Xsb_u+\delta rb_u=q_Xa+\delta v=0.
\]
因此 \([c_u]=[\delta]\,[r]\)，而 \([r][b_u]=[v]\)。这给出所写三角与标准三角的同构。

反方向，式\eqref{b4:eq:D-graph-inflation}是可容许短正合列。以 \(-\operatorname{pr}_{I_X}:Y\oplus I_X\to I_X\) 延拓 \(i_X\)，式\eqref{b4:eq:D-connecting-sign}在其余核上给出的连接态射就是 \(c_f\)。在稳定范畴中 \(Y\oplus I_X\cong Y\)，第一个映射成为 \([f]\)，第二个成为 \([b_f]\)，所以该短正合列给出标准三角。
\end{proof}

式\eqref{b4:eq:D-connecting-sign}中的负号配合本卷映射锥的末箭头约定。连接态射的稳定类记录扩张，而不是从原范畴任选一个 \(Z\to\Sigma X\) 的映射。

\subsection{三角结构及映射锥}
\label{b4:subsec:D03:02}

\begin{theorem}[Happel]\label{b4:thm:D-happel-triangulation}
设 \(\mathcal E\) 为局部小 Frobenius 范畴，固定式\eqref{b4:eq:D-syzygy-sequences}中的余合冲序列。以 \(\Sigma\) 为移位、以\cref{b4:def:D-stable-standard-triangle}中的好三角为指定三角，\(\underline{\mathcal E}\) 成为三角范畴。更换余合冲序列后，由比较映射给出的恒等对象函子连同 \(\Sigma\) 的比较自然同构是三角等价。
\end{theorem}

这一结果把正合范畴中的推出与本卷\cref{b4:def:triangulated-category}的三角公理联系起来。定理及复形例见 Bühler\cite[Section 13.4, Theorem 13.7 and Example 13.8]{Buehler2010ExactCategories}；该处的定理追溯到 Happel 专著第一章定理 2.6（第 16 页）的证明。

\begin{proofsketch}
\cref{b4:thm:D-stable-syzygy-equivalence}已经给出移位自等价。每个稳定态射由某个 \(f\) 表示，推出给出补全它的标准三角；\(1_X\) 的推出中 \(C_{1_X}\cong I_X\)，得到 \(X\xrightarrow{1}X\to0\to\Sigma X\)。若比较方块只在稳定范畴中交换，其误差通过一个投射内射对象分解；内射延拓使误差可以吸收到推出映射中，故可以构造三角之间的比较态射。旋转则通过把一条可容许短正合列与 \(Y\to I_Y\) 的序列共同推出，消去新增的投射内射项得到。

八面体所需的对象可以具体构造。对可复合的 \(X\xrightarrow{f}Y\xrightarrow{g}Z\)，令
\[
u=(f,-i_X):X\longrightarrow Y\oplus I_X,
\qquad
v(y,a)=(gy,-i_Yy,a):Y\oplus I_X\longrightarrow Z\oplus I_Y\oplus I_X.
\]
两者均为可容许单态射，\(v\) 是 \((g,-i_Y)\) 与 \(1_{I_X}\) 的直和，因此复合 \(vu\) 也可容许。令 \(T=\operatorname{coker}(vu)\)。对这两个可容许单态射及其复合应用正合范畴的 \(3\times3\) 引理，得到可容许短正合列
\[
0\longrightarrow\operatorname{coker}u\longrightarrow T
\longrightarrow\operatorname{coker}v\longrightarrow0.
\]
这一逐次商结论及 \(3\times3\) 引理分别见 Bühler\cite[Lemma 3.5; Corollary 3.6]{Buehler2010ExactCategories}。这里两端分别是 \(C_f\) 与 \(C_g\)。

为比较 \(T\) 与固定选择的 \(C_{gf}\)，利用 \(I_Y\) 的内射性，把 \(i_Yf:X\to I_Y\) 沿 \(i_X\) 延拓为 \(t:I_X\to I_Y\)，满足 \(ti_X=i_Yf\)。于是
\[
(i_Yf,i_X)=(t,1_{I_X})i_X:X\longrightarrow I_Y\oplus I_X.
\]
其中 \((t,1_{I_X})\) 与标准直和包含同构，所以这也是可容许单态射。复合 \(vu\) 的三个分量为 \((gf,-i_Yf,-i_X)\)。目标上的自同构
\[
\alpha(z,b,a)=(z,b-ta,a),
\qquad
\alpha^{-1}(z,b,a)=(z,b+ta,a)
\]
把它变为 \((gf,0,-i_X)\)，从而
\[
T\cong\operatorname{coker}(gf,0,-i_X)
\cong C_{gf}\oplus I_Y.
\]
因此 \(T\) 与 \(C_{gf}\) 在稳定范畴中同构。上面的可容许短正合列由\cref{b4:prop:D-conflation-triangle}给出八面体的第四个三角
\[
C_f\longrightarrow C_{gf}\longrightarrow C_g\longrightarrow\Sigma C_f.
\]
原来的两个包含与余核图同时提供其余三角间的映射。完整证明还须逐个核对旋转的负号、这些映射的交换性及更换嵌入后的相容性；这些公理核验采用所引 Happel 定理的证明，而以上构造说明八面体中的对象和短正合列从何而来。
\end{proofsketch}

\begin{proposition}\label{b4:prop:D-complex-triangle-identification}
设 \(\mathcal A\) 为局部小加性范畴，\(\operatorname{Ch}(\mathcal A)\) 配备逐项分裂正合结构，并取 \(I_X=\operatorname{Cone}(1_X)\)。则\cref{b4:thm:D-stable-complex-homotopy}中的识别
\[
\underline{\operatorname{Ch}(\mathcal A)}\cong K(\mathcal A)
\]
把 Frobenius 范畴的好三角与映射锥的好三角完全对应。
\end{proposition}
\begin{proof}
对 \(f:X\to Y\)，推出可以取 \(C_f=\operatorname{Cone}(f)\)。其两张结构映射逐次为
\[
Y\longrightarrow\operatorname{Cone}(f),\quad y\longmapsto(y,0),
\qquad
I_X\longrightarrow\operatorname{Cone}(f),\quad (x,x')\longmapsto(fx,x').
\]
它们在 \(X\) 上相同。给定从 \(Y,I_X\) 到复形 \(W\) 的相容映射 \(b,a\)，唯一的逐项推出映射为
\[
\operatorname{Cone}(f)^n\longrightarrow W^n,
\qquad (y,x')\longmapsto b^n(y)+a^n(0,x').
\]
在第二个分量上，其复形条件由 \(a\) 为复形映射及 \(a^n(x,0)=b^n f^n(x)\) 保证；第一个分量由 \(b\) 为复形映射保证。因此这确实是复形范畴中的推出。到 \(\Sigma X=X[1]\) 的映射为 \((y,x')\mapsto x'\)。所得标准三角正是本卷的映射锥三角。两类好三角均由这些标准三角的同构闭包定义，故完全相同。
\end{proof}

对逐项分裂短正合列，若选坐标 \(Y^n=X^n\oplus Z^n\)，写
\[
\mathrm{d}_Y^n=
\begin{pmatrix}\mathrm{d}_X^n&t^n\\0&\mathrm{d}_Z^n\end{pmatrix},
\]
则 \(a^n(x,z)=(x,t^nz):Y^n\to I_X^n\) 是延拓 \(i_X\) 的复形映射。式\eqref{b4:eq:D-connecting-sign}给出连接态射 \(\delta=-t:Z\to X[1]\)。它正是选分裂后由微分的非对角部分得到的映射，换分裂只改变其同伦类代表。

\ProseParagraphBreak
在双数代数中，取 \(i_k(1)=\varepsilon\)、\(I_k=A\) 及 \(\Sigma k=k\)。对式\eqref{b4:eq:D-dual-number-extension}可用 \(a=1_A\)，于是 \(\delta=-1_k\)。稳定范畴中的好三角为
\[
k\longrightarrow0\longrightarrow k\xrightarrow{-1_k}\Sigma k.
\]
因此中项变成零以后，非分裂扩张仍通过非零连接态射留下信息。稳定范畴既能表现模的扩张，也能表现复形的同伦；共同的原因是正合结构中的投射与内射对象可以同时消去。


\begin{chaptersummary}
稳定范畴保持对象，商掉经过投射内射对象的态射。核比较的差通过投射对象分解，余核比较的差通过内射对象分解；两类对象相同以后，合冲与余合冲成为互逆的加性自等价。正合结构仍是构造的一部分：同一对象范畴若只指定分裂列，所得稳定范畴可能完全为零。

\ProseParagraphBreak
双数代数的模由平方为零算子的二阶块和一阶块组成，自由的二阶块在稳定范畴中消失，一阶块之间仍有全部线性映射。对逐项分裂的复形，消失的对象恰为可缩复形，消失的态射恰为零同伦映射，所以稳定范畴就是同伦范畴。推出产生的好三角在这个模型中就是映射锥三角；在模模型中，它则保存了短正合扩张的连接资料。
\end{chaptersummary}

\BookExercises

\begin{exercise}\label{b4:ex:D-change-exact-structure}
设 \(k\) 为域，\(A=k[\varepsilon]/(\varepsilon^2)\)。在有限维左 \(A\)-模范畴上，分别取通常短正合结构和只含分裂列的结构。比较两者的投射内射对象及稳定范畴，并说明为什么不能只凭对象和态射决定稳定范畴。
\end{exercise}
\begin{solution}
通常结构中的投射内射对象恰为 \(A^r\)，而 \(k\) 不是投射对象；稳定范畴加性等价于有限维 \(k\)-向量空间，特别有 \(\operatorname{End}(k)=k\ne0\)。对分裂结构，任何态射沿分裂满射都能用截面提升，沿分裂单态射都能用收缩延拓，故所有对象均投射且内射。于是每个恒等映射都在稳定商中为零，所得范畴等价于零范畴。对象和原态射完全相同，指定的短正合列不同，消去的对象也就不同。
\end{solution}

\begin{exercise}\label{b4:ex:D-dual-number-block-maps}
设 \(k\) 为域，\(A=k[\varepsilon]/(\varepsilon^2)\)，\(M=A\oplus k\)，其中 \(k=A/(\varepsilon)\)。写出 \(M\) 的全部 \(A\)-线性自同态，计算普通自同态空间的维数和稳定自同态环。
\end{exercise}
\begin{solution}
以 \(A\oplus k\) 的顺序写分块。\(A\to A\) 是乘某个 \(a+b\varepsilon\)，\(k\to A\) 把 \(1\) 送到 \(c\varepsilon\)，\(A\to k\) 是商映射的 \(d\) 倍，\(k\to k\) 是乘 \(e\)。因此全部自同态可记为
\[
\begin{pmatrix}
a+b\varepsilon&c\varepsilon\\
d\pi&e
\end{pmatrix},
\qquad a,b,c,d,e\in k,
\]
其中 \(\pi:A\to k\) 是商映射。这五个参数独立，普通维数为五。凡定义域或陪域落在 \(A\) 分量的块，都通过投射模 \(A\) 分解；右下角的非零映射不能通过投射模分解，因为 \(k\to A^r\to k\) 总为零。故稳定自同态环为 \(k\)，对应参数 \(e\)。复合后的右下角是 \(ee'\)，因为交叉复合 \(k\to A\to k\) 为零，故这确为环同构。
\end{solution}

\begin{exercise}\label{b4:ex:D-no-enough-injectives}
给有限生成交换群范畴配备通常短正合结构。证明它有足够多投射对象，却没有足够多内射对象，因此不是 Frobenius 范畴。这里的内射性只针对该有限生成范畴中的短正合列检验。
\end{exercise}
\begin{solution}
每个有限生成交换群接收有限自由群的满射，而有限自由群投射，所以有足够多投射对象。若有限生成群 \(J\) 内射，对任意 \(n\ge1\) 和 \(x\in J\)，把 \(n\mathbb Z\to J\) 定义为 \(n\mapsto x\)。可容许单态射 \(n\mathbb Z\hookrightarrow\mathbb Z\) 的余核 \(\mathbb Z/n\) 仍有限生成，故该映射可延拓到 \(\mathbb Z\)，得到 \(y\in J\) 使 \(ny=x\)。于是 \(J\) 可除。

有限生成交换群写成 \(\mathbb Z^r\oplus T\)，其中 \(T\) 有限。若 \(r>0\)，乘二不是满射；若 \(T\ne0\)，乘 \(|T|\) 也不是满射。因此有限生成可除群只能为零，零对象是本范畴唯一内射对象。非零的 \(\mathbb Z\) 无法嵌入零对象，所以没有足够多内射对象。
\end{solution}

\begin{exercise}\label{b4:ex:D-truncated-cubic-period}
设 \(k\) 为域，\(A=k[t]/(t^3)\)，\(S=A/(t)\)。证明 \(A\) 是 Frobenius 代数，计算 \(\Omega S\)、\(\Omega^2S\) 及稳定自同态环；再证明 \(\Omega S\) 在稳定范畴中不与 \(S\) 同构。
\end{exercise}
\begin{solution}
令 \(\lambda\) 取 \(t^2\) 的系数，\(\phi(a)(b)=\lambda(ba)\)。若非零 \(a\) 的最低非零次数为 \(j\in\{0,1,2\}\)，乘 \(t^{2-j}\) 后的 \(t^2\) 系数非零，故 \(\phi\) 单射；维数相同使它同构，且左线性，证明 Frobenius 性。

用 \(A\to S\) 的商映射，得到 \(\Omega S=tA\)。乘 \(t\) 的满射 \(A\to tA\) 的核为 \(t^2A\cong S\)，故 \(\Omega^2S\cong S\)。任何 \(S\to A\) 的像落入 \(t^2A\)，任何 \(A\to S\) 都把它送到零，因此稳定自同态环为 \(k\)。投射模是有限自由模的直和项，所以经过任意投射模的 \(S\to S\) 映射也为零。

任意 \(S\to tA\) 的像都在 \(t^2A\)，任意 \(tA\to S\) 都杀死 \(t^2A=t(tA)\)，故它们的复合为零。若二者稳定同构，就应有这样的两映射，其复合的稳定类为 \(1_S\)；但 \([1_S]\ne0\)，矛盾。这里合冲的二次迭代回到 \(S\)，一次迭代则没有回到它。
\end{solution}

\begin{exercise}\label{b4:ex:D-same-graded-different-contractibility}
设 \(k\) 为域。考虑仅在次数零和一次非零、两项都为 \(k\) 的上链复形，分别令微分为 \(1_k\) 和 \(0\)。计算它们在逐项分裂复形范畴的稳定范畴中是否为零。
\end{exercise}
\begin{solution}
微分为 \(1_k\) 时，令 \(h^1=1_k\)，其余收缩分量为零，得到 \(\mathrm{d}h+h\mathrm{d}=1\)，所以可缩，在稳定范畴中为零。微分为零时，任意 \(h\) 都有 \(\mathrm{d}h+h\mathrm{d}=0\)，不可能等于非零的恒等映射，故不可缩，在稳定范畴中不为零。两者的分次对象相同，但微分决定了完全不同的稳定对象。
\end{solution}

\begin{exercise}\label{b4:ex:D-explicit-nullhomotopy-factorization}
令 \(X\) 为交换群上链复形 \(X^0=X^1=\mathbb Z\)、\(\mathrm{d}_X^0=2\)，其余项为零。令 \(f:X\to X\) 在两项上都是乘二。写出 \(f\) 的零同伦，以及经过 \(I_X=\operatorname{Cone}(1_X)\) 的具体分解，并核对其中映射与微分相容。
\end{exercise}
\begin{solution}
取 \(h^1:X^1\to X^0\) 为恒等，其余 \(h^n=0\)。零次有 \(h^1\mathrm{d}_X^0=2\)，一次有 \(\mathrm{d}_X^0h^1=2\)，故 \(f=\mathrm{d}h+h\mathrm{d}\)。可缩复形 \(I_X\) 的非零部分为
\[
\mathbb Z\xrightarrow{x\mapsto(x,-2x)}\mathbb Z^2
\xrightarrow{(a,b)\mapsto2a+b}\mathbb Z,
\]
分别在次数 \(-1,0,1\)。包含 \(i_X\) 在零次把 \(x\) 送到 \((x,0)\)，在一次是恒等。定义 \(b:I_X\to X\) 的零次分量为 \(b^0(a,b')=2a+b'\)，一次分量为 \(b^1(c)=2c\)，负一次分量为零。零次分量作用于前一个微分得到 \(2x-2x=0\)；后一方块的两条路径都为 \(2(2a+b')\)，所以 \(b\) 是复形映射。两项上的 \(bi_X\) 都为乘二，给出所求分解。
\end{solution}

\begin{exercise}\label{b4:ex:D-acyclic-not-stably-zero}
设 \(k\) 为域，\(A=k[\varepsilon]/(\varepsilon^2)\)。令双向无界复形 \(X\) 在每个整数次数都是左正则模 \(A\)，每个微分都是乘 \(\varepsilon\)。证明 \(X\) 零调但不可缩，并解释它在同伦范畴与导出范畴中的不同表现。
\end{exercise}
\begin{solution}
\(\varepsilon^2=0\) 保证复形条件。每个次数的核与像都等于 \(\varepsilon A\)，所以所有上同调为零。若有 \(A\)-线性收缩，每个 \(h^n:A\to A\) 都是乘某个 \(a_n\in A\)。收缩等式要求
\[
\varepsilon(a_n+a_{n+1})=1
\]
在每个次数成立，而左边在理想 \((\varepsilon)\) 内，右边不在，矛盾。因此 \(X\) 在 \(K(A\text{-}\mathbf{Mod})\) 及相应逐项分裂稳定范畴中不为零。到零复形的映射却是拟同构，所以 \(X\) 在 \(D(A\text{-}\mathbf{Mod})\) 中为零。稳定复形商消去的是可缩复形，导出局部化还会消去这个不可缩的零调复形。
\end{solution}

\begin{exercise}\label{b4:ex:D-extension-parameter-connecting}
设 \(k\) 为域，\(A=k[\varepsilon]/(\varepsilon^2)\)。对 \(a\in k\)，令 \(E_a=ku\oplus kv\)，规定 \(\varepsilon u=0\)、\(\varepsilon v=au\)。把它置于扩张
\[
0\longrightarrow k\xrightarrow{1\mapsto u}E_a
\xrightarrow{u\mapsto0,\ v\mapsto1}k\longrightarrow0.
\]
证明这些参数给出 \(\operatorname{Ext}_A^1(k,k)\cong k\)，确定何时分裂，并在 \(i_k(1)=\varepsilon\)、\(\Sigma k=k\) 的选择下求连接态射。
\end{exercise}
\begin{solution}
任何这样的扩张在向量空间上都分裂，选取左端基元的像 \(u\)，再选 \(1\) 的提升 \(v\)。因商模的 \(\varepsilon\) 作用为零，必有 \(\varepsilon v=au\)；更换提升 \(v\mapsto v+cu\) 不改变 \(a\)。固定两端的扩张同构必须把 \(u\) 送到 \(u\)，把 \(v\) 送到 \(v+cu\)，与 \(\varepsilon\) 作用相容便要求两参数相同。因此全部扩张类恰由 \(a\in k\) 参数化。

Baer 和先取纤维积，再沿两份核的加法推出。纤维积有核基元 \(u_1,u_2\) 及商提升 \((v_a,v_b)\)，其 \(\varepsilon\) 像为 \(au_1+bu_2\)；推出后两个核基元都成为同一个 \(u\)，所以新参数为 \(a+b\)。沿核的标量 \(c\) 推出则把参数变为 \(ca\)。这给出所称 \(k\)-线性同构。\(a=0\) 时 \(E_a=k\oplus k\)，列分裂；\(a\ne0\) 时由 \(1\mapsto v\)、\(\varepsilon\mapsto au\) 得 \(A\cong E_a\)，列不分裂。

延拓 \(i_k\) 可取 \(\alpha:E_a\to A\)，\(\alpha(u)=\varepsilon\)、\(\alpha(v)=a\)。它左线性，因为 \(\alpha(\varepsilon v)=a\varepsilon=\varepsilon\alpha(v)\)。商映射 \(q:A\to k\) 使 \(q\alpha(v)=a\)，故式\eqref{b4:eq:D-connecting-sign}给出 \(\delta=-a:k\to\Sigma k\)。稳定 Hom 为 \(k\)，所以非零扩张类恰好给出非零连接态射。
\end{solution}
